Popper and Hermann may have had another opportunity to meet at the Second International Congress for the Unity of Science, held in Copenhagen in June 1936, where \citet{Hermann1936} delivered a short address. However, there is no evidence that the encounter left any trace. They both left continental Europe shortly thereafter, in 1937. They revisited \qm only after the war, around the turn of the 1950s \citep{Henry-Hermann1948,Popper1951}. Whereas Hermann seems to have moved on, absorbed in political and pedagogical commitments, for Popper this marked the beginning of a full-scale return to the philosophy of \qm after the debacle of the 1930s \citep{DelSanto2019}. Popper began writing the material on \qm in 1951--56, shortly before the publication of the first English edition of the \Lo in 1959. Part of this material was included in the appendices, while the rest became a separate \textit{Postscript} to the \Lo, which reached galley proofs in 1956--57 but remained unpublished. However, 

%%Popper never understood the deep of Hermann's critique to the notion of trajectory, that the classical notion of trajectory breaks down because quantum-mechanical determinations belonging to different experimental contexts cannot be combined. 
%
%The reason is probably more clearly by Pauli’s later argument brings out a further consequence of this incompatibility: unlike classical measurements, quantum measurements are not cumulative. A new measurement does not merely add information about the same trajectory but establishes a new predictive situation, potentially rendering information obtained from earlier measurements unusable. Even in the most 

%Hermann does not formulate the point in these terms. Nearly two decades later, however, Pauli drew a particularly illuminating contrast between classical and quantum mechanics. In classical mechanics, repeated measurements of finite accuracy can be combined to determine a trajectory with progressively greater accuracy. In quantum mechanics, by contrast, each new position measurement introduces a new momentum uncertainty and can thereby destroy the usefulness of earlier measurements for subsequent predictions. Pauli thus makes explicit the non-cumulative character of quantum-mechanical measurements that is implicit in Hermann’s rejection of a context-independent trajectory.

%The strongest defensible formulation may be “provides a later perspective from which the significance of Hermann’s argument becomes particularly clear.

In the new English edition of the \Lo, the main text of the chapter on \qm was left untouched; however, Popper sought to make amends for his previous mistakes with new footnotes and appendices marked with starred numbers. Popper acknowledged what \vW had pointed out to him: trajectory reconstruction was not possible \q{\myemph{before} the first of these measurements}, that is, the momentum measurement \citep[242\fnstar{3}\me]{Popper1959}. Without this, Popper's \TE fails. However, Popper continued to hold that \q{non-predictive measurements determine the path of a particle only \myemph{between} two measurements,} such as a measurement of momentum followed by one of position \citep[242\fnstar{3}\me]{Popper1959}. Popper seems not to have appreciated, and quite possibly not to have read, Hermann's deeper critique that quantum mechanics does not license trajectory reconstructions in general. This is not simply an oversight. Popper’s interpretation of \qm had changed significantly after the war. The frequency interpretation of probability had been replaced by the propensity interpretation \citep{Popper1957,Popper1959}, the issue of determinism by that of realism, and the language of particles by that of experimental arrangements \citep{Howard2012}. However, Popper never ceased to believe that Heisenberg’s concession that the existence of past trajectories was a \s{matter of taste} opened a crack in the \s{orthodox interpretation} that Popper hoped to widen.

The 1956--57 draft of the \emph{Postscript}, and even more clearly a 1966 article \citep{Popper1967} conceded that in \qm \emph{predictions} of future trajectories are statistical, with dispersions related by the Heisenberg formulae. However, testing these predictions requires measurements, and \q{\emph{these measurements are retrodictions}} (\cite[25]{Popper1967}; \cite[60]{Popper1982}); such retrodictions are, in Popper's view, reconstructions of particle trajectories. Since \qm is a theory concerned with the probabilities of trajectories, trajectories must be real; \ie, the particles therefore \s{possess} sharp positions and momenta at every instant. While the \UR forbid combining incompatible arrangements to \emph{prepare} an ensemble of particles with the same position and momentum without producing a scatter, Popper insisted that it must be possible to combine the corresponding arrangements to \emph{measure} the very scatter predicted by the \UR \citep[21]{Popper1967}. According to \citet[III-17]{Popper1982}, this is precisely what the EPR \TE demonstrates. However, if, upon freely choosing to measure either the position or the momentum of system $A$, one can determine with certainty, in two different runs, both the position and the momentum of $B$ \emph{without disturbing} it, then $B$ possesses definite values for position and momentum, even though the theory does not enable us to predict them simultaneously.

%According to Popper, however, EPR had already conceded this premise: 


In Popper's view, \citets[697]{Bohr1935} objection surreptitiously changes the argument. Even in the absence of a \s{mechanical} disturbance, he takes Bohr to argue, the counterfactual status of the unperformed measurement constitutes a form of \s{non-mechanical disturbance} that prevents us from inferring simultaneous elements of reality \citep[445]{Popper1959}. If we choose to measure the momentum of $A$, we can determine the momentum of $B$ with arbitrary precision without mechanically disturbing it; however, in doing so we forgo the possibility of measuring its position. Had we instead chosen to measure position, the opposite would have been the case: \q{The arrangements are exclusive of each other (or \s{complementary}), the argument continues, and so we cannot combine the two spaces \textelp{} into one} \citep[149]{Popper1982}. Popper, however, regarded this prohibition as contrary to common sense \citep[150]{Popper1982}: if I can freely choose to determine either your weight or your height with certainty without disturbing you, then it is unreasonable to deny that you possess both simultaneously \citep[150]{Popper1982}. Popper's choice of a counterexample is somewhat careless: height and weight can be measured within the same experimental arrangement in a single run, whereas in \qm, position and momentum cannot. Nevertheless, it reveals that Popper had come to realize that the question of the reality of past trajectories ultimately boils down to the relationship between experimental arrangements \citep[54]{Howard2012}.

In retrospect, this issue lay behind the Hermann-Popper debate all along. As Popper put it, testing a statistical theory requires \s{\emph{retrodictions}}. He equated retrodictions with reconstructions of particle \s{\emph{trajectories}}: trajectories must be real, and particles must therefore possess simultaneous position and momentum \emph{despite} the mutual exclusivity of the corresponding experimental arrangements. Hermann identified the questionable premise on which this argument rests. Retrodictions are indeed indispensable for testing \qm, but they reconstruct \s{\emph{measurement}} processes: \emph{because} the corresponding experimental arrangements are mutually exclusive, particles cannot possess simultaneous position and momentum, and trajectories are therefore meaningless. Although the controversy over trajectories later largely receded\footnote{See, however, the criticism of Popper in \cite{Feyerabend1968,Feyerabend1969} and his defense in \cite[366--67]{Ballentine1970}}, it foreshadowed the deeper philosophical issue that emerged after EPR. Popper came to realize that the deeper source of his unease with quantum mechanics lay not in its challenge to determinism but in its challenge to realism. He shared Einstein's fear that, unless the physical state could be fully detached from the experimenter's intervention, one would have to accept the \q{intrusion of subjectivism into physics} \citep[2, 5]{Popper1982}. Hermann's great contribution to the early philosophical reception of \qm was not so much her somewhat contrived defense of causality; it was her keen recognition that quantum mechanics requires a weaker form of detachment. State attribution depends on the experimenter's free choice among mutually exclusive experimental arrangements, yet this does not render the physical state \s{subjective}. Once the experimental arrangement is fixed, the state attribution is objectively determined, albeit \s{relative to the observational context} \citep[380--81]{Henry-Hermann1948}.

\nocite{Popper2006}



%Unabhangig von Einstein noch eine Bemerkung zur Illustration des Unter-
%schiedes von klassischer Mechanik und Quantenmechanik bei der "Messung"
%einer "Bahn", die Sie in Brief und Manuskript auch behandeln.
%A. Klassische Mechanik
%Denken wir z. B. an die Bestimmung der Bahn eines Planeten: Man messe
%nun wiederholt (in verschiedenen Zeitmomenten to, tl...) den Ort immer wieder
%mit derselben Genauigkeit L1xo. Hat man ein
%einfaches Gesetz rur die Bewegung des Korpers (z. B. das Newtonsche Gravita-
%tionsgesetz), so kann man daraus die Bahn (also Lagen und Geschwindigkeiten
%zu irgendeiner Zeit) des Korpers mit beliebig hoher Genauigkeit berechnen (und
%das angenommene Gesetz zu spateren oder friiheren Zeiten als die beniitzten
%auch wieder priifen). Wiederholte Ortsmessungen mit begrenzter  kannen daher hier eine Ortsmessung mit hoher Genauigkeit erfolgreich ersetzen.
%Die Annahme relativ einfacher Kraftgesetze wie des Newtonschen (und nicht
%irgendwelche irreguHiren Zickzackbewegungen im Kleinen) erscheint dabei als
%eine im Sinne der klassischen Mechanik erlaubte Idealisation.
%B. Quantenmechanik
%Die Wiederholung der Ortsmessung der aufeinanderfolgenden Zeiten mit
%derselben Genauigkeit L\xo niitzt gar nichts fur die Vorhersagbarkeit spaterer
%Ortsmessungen. Dennjede Ortsmessung mit Genauigkeit L\xo zur Zeit tn bringt
%die Ungenauigkeit9
%h
%L\xn+lrv --(tn+l- tn )
%mL\xo
%zur spateren Zeit wieder mit sich und vernichtet die Verwertbarkeit aZZerfriiheren
%Ortsmessungen innerhaZb dieser FehZergrenze. (Wenn ich nicht irre, hat dieses
%Beispiel vor vielen Jahren schon Bohr mit mir diskutiert.)
%Der Hauptunterschied beider Theorien A und B, daB namlich bei B durch eine
%Messung Kenntnisse auf Grund friiherer Messungen verlorengehen kannen, ist
%in Ihrem Manuskript nicht geniigend zum Ausdruck gekommen!
%Viele GriiBe Stets Ihr W. PauliGenauigkeit

